\documentclass[UTF8,a4paper,11pt,fontset=fandol]{ctexart}
\usepackage[margin=24mm]{geometry}
\usepackage{amsmath,amssymb,booktabs,tabularx,array}
\usepackage[unicode,hidelinks]{hyperref}
\pagestyle{plain}
\linespread{1.10}
\setlength{\parskip}{0.2em}
\setlength{\emergencystretch}{2em}
\renewcommand{\arraystretch}{1.12}
\newcolumntype{Y}{>{\raggedright\arraybackslash}X}
\newcommand{\kw}[1]{\ensuremath{\mathsf{#1}}}
\newcommand{\infer}[2]{\displaystyle\frac{\displaystyle #1}{\displaystyle #2}}
\usepackage{xcolor}
\usepackage{tikz}
\usetikzlibrary{arrows.meta,calc,positioning}
\definecolor{gmFill}{HTML}{FFF8E6}
\definecolor{gmBorder}{HTML}{B7791F}
\definecolor{aicFill}{HTML}{EEF5FC}
\definecolor{aicMem}{HTML}{DCEBFA}
\definecolor{aicBorder}{HTML}{2B6CB0}
\definecolor{aivFill}{HTML}{EDF9F4}
\definecolor{aivMem}{HTML}{D4F0E2}
\definecolor{aivBorder}{HTML}{2F855A}
\definecolor{scalarFill}{HTML}{EDF2F7}
\definecolor{scalarBorder}{HTML}{4A5568}
\definecolor{syncLocal}{HTML}{1D4ED8}
\definecolor{syncBar}{HTML}{C53030}
\definecolor{syncCross}{HTML}{6B46C1}
\definecolor{coreBg}{HTML}{F8FAFC}
\definecolor{coreBorder}{HTML}{94A3B8}
%%%%%%%%%%%%%% Notation macros %%%%%%%%%%%%%%
\newcommand{\an}[3]{{\color{#2} \textsc{#1}:#3}}
\newcommand{\litao}[1]{\an{litao}{red}{#1}}

\hypersetup{pdftitle={Core Language Operational Semantics},pdfauthor={}}
\begin{document}
\begin{center}
\Large\bfseries 核心语言的操作语义
\end{center}


\noindent

A \emph{core language} for describing the behavior of a single \emph{AI Core}.

\[
\begin{aligned}
1\ \mathrm{NPU}\ (\mathrm{rank})
  &= 24\ \text{AI Cores}\\
  &= 24\times(\mathrm{AIC}+2\,\mathrm{AIV}).
\end{aligned}
\]
Each AIC/AIV core has several pipes, for example
$
\mathsf{PIPE\_MTE2},\ \mathsf{PIPE\_V},\ \mathsf{PIPE\_MTE3},\ \ldots
$

This document describes a core language slightly modified from the initial
version in the \texttt{docs/core-language.tex} file, and presents an operational semantics.

\section{参数说明}

对单个\textbf{硬件配置}来说，可配置的参数如下。
右列以 Ascend 910B 实际取值为例。
本文定义的良构检查和操作语义均基于上述参数，不固定一代硬件具体参数的取值。
记完整配置为 $\mathcal H$：它包含本节的硬件参数，以及后文的静态签名、允许的操作/同步关系和动态解释。
类型判断写作 $\mathcal H;\Gamma\vdash_s C$，小步关系写作 $\longrightarrow_{\mathcal H}$；
未加下标的 $\mathsf{kind}$、$\mathsf{Pipe}$ 等均取自当前 $\mathcal H$。

$f$ 是发到数据管道上的运算，由该管道异步完成。$g$ 是标量算子，由程序直接执
行，\kw{let} 立刻绑定结果。二者都可根据IR提供的实际运算配置。
本文把同步原语（\kw{setflag}、\kw{waitflag}、\kw{pipebarrier}、\kw{syncblockset} 和
\kw{syncblockwait}）作为核心语言的原语。
\begin{equation}
\label{eq:hw}
\begin{array}{@{}lll@{}}
\text{参数说明} &  & \text{Example: Ascend 910B}\\
\hline
\text{Subcore 种类}
  & d\in\mathsf{CoreKind}
  & \{\mathsf{AIC},\mathsf{AIV}\}\\[0.3em]
\text{硬件实际的指令流}
  & s\in\mathsf{Stream}
  & \{\mathsf{aic},\mathsf{aiv}_1,\mathsf{aiv}_2\}\\[0.3em]
\text{指令流到Subcore映射 }
  & \mathsf{kind}:\mathsf{Stream}\to\mathsf{CoreKind}
  & \begin{aligned}
      \mathsf{aic}&\mapsto\mathsf{AIC}\\
      \mathsf{aiv}_1,\mathsf{aiv}_2&\mapsto\mathsf{AIV}
    \end{aligned}\\[0.6em]
\text{直接执行的标量算子}
  & g\in\mathsf{Scalar}
  & \begin{aligned}
      &\{\texttt{arith},\texttt{get\_block\_idx},\\
      &\quad\texttt{get\_sub\_block\_idx},\ldots\}
    \end{aligned}\\[0.6em]
\text{发到管道的运算}
  & f\in\mathsf{Operator}_d
  & \begin{aligned}
      \mathsf{Operator}_{\mathsf{AIV}}
        &=\{\mathsf{Vadd},\mathsf{Load},\mathsf{Store},\ldots\}\\
      \mathsf{Operator}_{\mathsf{AIC}}
        &=\{\mathsf{Mmad},\mathsf{Fixpipe},\ldots\}
    \end{aligned}\\[0.6em]
\text{SubCore的地址空间}
  & \mathsf{Space}_d
  & \begin{aligned}
      \mathsf{Space}_{\mathsf{AIV}}
        &=\{\mathsf{GM},\mathsf{UB}\}\\
      \mathsf{Space}_{\mathsf{AIC}}
        &=\{\mathsf{GM},\mathsf{L1},\mathsf{L0A},\mathsf{L0B},\mathsf{L0C}\}
    \end{aligned}\\[0.6em]
\text{跨流共享地址空间}
  & \mathsf{SharedSpace}_{\mathcal H}
  & \{\mathsf{GM}\}\\[0.3em]
\text{数据管道}
  & \mathsf{Pipe}_d
  & \begin{aligned}
      \mathsf{Pipe}_{\mathsf{AIV}}
        &=\{\mathsf{MTE2},\mathsf{V},\mathsf{MTE3}\}\\
      \mathsf{Pipe}_{\mathsf{AIC}}
        &=\{\mathsf{MTE1},\mathsf{MTE2},\mathsf{M},\mathsf{FIX}\}
    \end{aligned}\\[0.6em]
\text{同步管道名}
  & \mathsf{SyncName}
  & \{\mathsf{S},\mathsf{ALL}\}
\end{array}
\end{equation}

管道名包括数据管道和同步管道，我们统一使用元变量$p,q$ 表示如下：
\begin{equation}
\label{eq:pipe}
p,q\in\mathsf{Pipe}
  =\Bigl(\bigcup_{d\in\mathsf{CoreKind}}\mathsf{Pipe}_d\Bigr)
   \cup\mathsf{SyncName}.
\end{equation}



对单个\textbf{程序}来说，指定它唤起的指令流 $\kappa\subseteq\mathsf{Stream}$。程序不重新指定种类，一条流的种类是 $\mathsf{kind}(s)$。HIVM 的 \texttt{core\_ratio} 记录过的比例包括 $0{:}1$、$1{:}0$、$1{:}1$ 和 $1{:}2$，分别对应
\[
\begin{aligned}
\kappa_{0:1}&=\{\mathsf{aiv}_1\},&
\kappa_{1:0}&=\{\mathsf{aic}\},\\
\kappa_{1:1}&=\{\mathsf{aic},\mathsf{aiv}_1\},&
\kappa_{1:2}&=\{\mathsf{aic},\mathsf{aiv}_1,\mathsf{aiv}_2\}.
\end{aligned}
\]
纯 AIV 或纯 AIC 的 kernel 使用 $\kappa_{0:1}$ 或 $\kappa_{1:0}$。后文的语法和规
则对任意 $\kappa\subseteq\mathsf{Stream}$ 相同。一项程序是 $\kappa$ 中这些指令流
的并列，不包含第二个 AI Core。

同一个 AI Core 里有两层同步。
\kw{setflag}、\kw{waitflag} 与 \kw{pipebarrier} 只约束一条指令流内部的管道。
\kw{syncblockset} 与 \kw{syncblockwait} 约束这个 AI Core 内的 AIC 与 AIV。
Ascend C 把这一层也叫做核间同步，对应 \texttt{CrossCore} 的模式 2。
本文保留这两层。
多个 AI Core 之间的同步不在本文中，那是 \texttt{CrossCore} 的模式 0，以及 block 网格和 rank 之间的通信。

目前的局限：
\begin{itemize}
\item 配置只含一个 AI Core。\texttt{get\_block\_idx} 是这个 AI Core 已经给定的整数，各条流读到同一个值。\texttt{get\_sub\_block\_idx} 是一条流自己的整数，用来区分 $\kappa$ 里的各条 AIV 流。
\item 不处理复杂算子如 \texttt{mmadL1} 的内部事件槽、\texttt{nd2nz} 和 FFTS 基址。写到 $\mathsf{GM}$ 的 \texttt{fixpipe} 是 $\mathsf{Operator}_{\mathsf{AIC}}$ 中的一个名字，仍是一条 \kw{issue}。
\item 标量不直接读写共享内存 $\mu$，因此核内以 $\mathsf{S}$ 为端点的 \texttt{set\_flag} 与 \texttt{wait\_flag}（由 \texttt{memref.load}、\texttt{memref.store} 触发）不纳入基础配置，可作为扩展加入。
\end{itemize}

\section{语法}

沿用 \eqref{eq:hw}、\eqref{eq:pipe} 中定义的参数。
\[f\in\mathsf{Operator}_d
\quad g\in\mathsf{Scalar}
\quad m\in\mathsf{Space}_d
\quad p,q\in\mathsf{Pipe}
\quad \text{where } d = \mathsf{kind}(s)\ (s \in \kappa)
\]
\[
\begin{aligned}
\tau ::={}
  &\ \mathsf{int}\mid\mathsf{float}\mid\mathsf{bool}\mid\mathsf{event}\mid\mathsf{view}\\
v ::={}
  &\ c\mid x\\
C ::={}
  &\ \kw{return}(\bar v)
   \mid \kw{end}\\
  &\mid \kw{let}\ x=\kw{scalar}[g](\bar v)\ \kw{in}\ C
   \mid \kw{let}\ x=\kw{alloc}[m](\bar v)\ \kw{in}\ C\\
  &\mid \kw{issue}[p](f,\bar v);\ C\\
  &\mid \kw{setflag}[p,q](v);\ C
   \mid \kw{waitflag}[p,q](v);\ C\\
  &\mid \kw{pipebarrier}[p];\ C\\
  &\mid \kw{syncblockset}[p,q](v);\ C
   \mid \kw{syncblockwait}[p,q](v);\ C\\
  &\mid \kw{if}\ v\ \kw{then}\ \{C_1\}\ \kw{else}\ \{C_2\};\ C_3\\
  &\mid \kw{for}\ x=v_1\ \kw{to}\ v_2\ \kw{step}\ v_3\ \kw{do}\ \{C_1\};\ C_2,\\[0.3em]
P ::={}
  &\ \mathop{\Big\|}_{s\in\kappa} s\{C_s\}.
\end{aligned}
\]

\begin{itemize}
  \item $P$ 并行执行若干条指令流，$\kappa$ 只有一个元素时，$P$ 就写作
  $s\{C\}$。每条指令流内顺序执行。
  \item $g$ 的例子包括：标量算术运算、\texttt{get\_block\_idx}、\texttt{get\_sub\_block\_idx} 和视图构造。$\kw{scalar}[g]$ 与 $\kw{alloc}[m]$ 均不进入数据管道，由 $\kw{let}$ 立刻绑定结果；$\kw{alloc}[m](\bar v)$ 在空间 $m$ 中分配一块新鲜存储并绑定为 $\mathsf{view}$。
  \item $\kw{issue}[p](f,\bar v)$ 不绑定结果。$\bar v$ 里的地址是类别为
  $\mathsf{view}$ (可理解为 \texttt{memref} 等地址表示) 的操作数，读取和写入都在
  其中，目的地也是一个操作数。后续使用 $\mathsf{foot}(s,p,f,\bar\nu)=(R,W)$ 计算视图
  对应的存储单元集合（表示为 $\mathsf{Loc}$ 的有限子集）。令 $\mathsf{Space}=\bigcup_d\mathsf{Space}_d$，
  按跨流共享空间与流私有空间区分具体地址：
  \[
  \begin{aligned}
  \mathsf{Loc}_{s,m} &=
    \begin{cases}
      \{(m,z)\mid z\in\mathbb Z\}, & m\in\mathsf{Space}_{\mathsf{kind}(s)}\cap\mathsf{SharedSpace}_{\mathcal H},\\
      \{(s,m,z)\mid z\in\mathbb Z\}, & m\in\mathsf{Space}_{\mathsf{kind}(s)}\setminus\mathsf{SharedSpace}_{\mathcal H},
    \end{cases}\\
  \mathsf{Loc}_s &= \bigcup_{m\in\mathsf{Space}_{\mathsf{kind}(s)}}\mathsf{Loc}_{s,m},
  \qquad
  \mathsf{Loc} = \bigcup_{s\in\mathsf{Stream}}\mathsf{Loc}_s,
  \qquad
  R,W\subseteq\mathsf{Loc}_s.
  \end{aligned}
  \]
  其中 $R$ 是读取的单元，$W$ 是写入的单元。$\mathsf{Loc}$ 只出现在操作语义里；同一 AI Core 中不同 AIV 流的私有 $\mathsf{UB}$ 由流标识 $s$ 隔离，而 $\mathsf{GM}$ 等共享空间由各流共享。
\end{itemize}

\litao{We may further abstract the event semantics, using a more general event
identifier, while for Ascend instantiate it to $(p,q,\varepsilon)$.}


% \clearpage
\section{良构与基本类型检查}
\label{sec:wf}
本节检查基本类别、变量作用域、计算单元与管道归属、\kw{alloc} 的地址空间归属，以及配置允许的操作和同步关系；不检查 view 的具体布局或地址边界。
固定合法配置 $\mathcal H$ 与有限活动流集合 $\kappa\subseteq\mathsf{Stream}$。
判断 $\mathcal H;\Gamma\vdash_s C$ 检查指令流 $s$ 中命令 $C$ 的类别与作用域；规则使用的配置参数及同步条件在下表中解释。
\begingroup\small
\setlength{\abovedisplayskip}{5pt}\setlength{\belowdisplayskip}{5pt}
\begin{equation*}
\begin{gathered}
\infer{\Gamma(x)=\tau}{\mathcal H;\Gamma\vdash x:\tau}\;\hypertarget{wf-var}{\textsc{Var}}
\qquad
\infer{\mathsf{lit}_{\mathcal H}(c)=\tau}{\mathcal H;\Gamma\vdash c:\tau}\;\hypertarget{wf-lit}{\textsc{Lit}}
\qquad
\infer{|\bar v|=|\bar\tau|\quad\forall i.\ \mathcal H;\Gamma\vdash v_i:\tau_i}
 {\mathcal H;\Gamma\vdash\bar v:\bar\tau}\;\hypertarget{wf-tuple}{\textsc{Tuple}}\\[1em]
\infer{\Delta_{\mathcal H}(g)=(\bar\tau,\tau)\quad\mathcal H;\Gamma\vdash\bar v:\bar\tau\quad
 x\notin\mathrm{dom}(\Gamma)\quad\mathcal H;\Gamma,x:\tau\vdash_s C}
 {\mathcal H;\Gamma\vdash_s\kw{let}\ x=\kw{scalar}[g](\bar v)\ \kw{in}\ C}
 \;\hypertarget{wf-let}{\textsc{Let}}\\[1em]
\infer{m\in\mathsf{Space}_{\mathsf{kind}(s)}\quad\mathcal H;\Gamma\vdash\bar v:\mathsf{int}^*\quad
 x\notin\mathrm{dom}(\Gamma)\quad\mathcal H;\Gamma,x:\mathsf{view}\vdash_s C}
 {\mathcal H;\Gamma\vdash_s\kw{let}\ x=\kw{alloc}[m](\bar v)\ \kw{in}\ C}
 \;\hypertarget{wf-alloc}{\textsc{Alloc}}\\[1em]
\infer{p\in\mathsf{Pipe}_{\mathsf{kind}(s)}\quad f\in\mathsf{Operator}_{\mathsf{kind}(s)}\quad
 \Sigma_{\mathcal H}(\mathsf{kind}(s),p,f)=\bar\tau\quad\mathcal H;\Gamma\vdash\bar v:\bar\tau\quad\mathcal H;\Gamma\vdash_s C}
 {\mathcal H;\Gamma\vdash_s\kw{issue}[p](f,\bar v);\ C}
 \;\hypertarget{wf-issue}{\textsc{Issue}}\\[1em]
\infer{\mathsf{ok}_{\mathcal H,s}(a)\quad\mathcal H;\Gamma\vdash v:\mathsf{event}\quad\mathcal H;\Gamma\vdash_s C}
 {\mathcal H;\Gamma\vdash_s a;\ C}\;\hypertarget{wf-sync}{\textsc{Sync}}
\qquad
\infer{(\mathsf{kind}(s),p)\in\mathsf{Barrier}_{\mathcal H}\quad\mathcal H;\Gamma\vdash_s C}
 {\mathcal H;\Gamma\vdash_s\kw{pipebarrier}[p];\ C}\;\hypertarget{wf-bar}{\textsc{Bar}}\\[1em]
\infer{\mathcal H;\Gamma\vdash v:\mathsf{bool}\quad\mathcal H;\Gamma\vdash_s C_1\quad\mathcal H;\Gamma\vdash_s C_2\quad\mathcal H;\Gamma\vdash_s C_3}
 {\mathcal H;\Gamma\vdash_s\kw{if}\ v\ \kw{then}\ \{C_1\}\ \kw{else}\ \{C_2\};\ C_3}
 \;\hypertarget{wf-if}{\textsc{If}}\\[1em]
\infer{\mathcal H;\Gamma\vdash v_1,v_2,v_3:\mathsf{int},\mathsf{int},\mathsf{int}\quad
 x\notin\mathrm{dom}(\Gamma)\quad\mathcal H;\Gamma,x:\mathsf{int}\vdash_s C_1\quad\mathcal H;\Gamma\vdash_s C_2}
 {\mathcal H;\Gamma\vdash_s\kw{for}\ x=v_1\ \kw{to}\ v_2\ \kw{step}\ v_3\ \kw{do}\ \{C_1\};\ C_2}
 \;\hypertarget{wf-for}{\textsc{For}}\\[1em]
\infer{\mathcal H;\Gamma\vdash\bar v:\bar\tau}{\mathcal H;\Gamma\vdash_s\kw{return}(\bar v)}
 \;\hypertarget{wf-return}{\textsc{Return}}
\qquad
\infer{}{\mathcal H;\Gamma\vdash_s\kw{end}}\;\hypertarget{wf-end}{\textsc{End}}\\[1em]
\infer{\forall s\in\kappa.\ \mathcal H;\varnothing\vdash_s C_s}
 {\mathcal H\vdash_\kappa\mathop{\Big\|}_{s\in\kappa}s\{C_s\}}
 \;\hypertarget{wf-program}{\textsc{Program}}
\end{gathered}
\end{equation*}
\endgroup

\noindent 规则中的环境与配置参数如下。$\mathsf{Type}$ 是语法给出的五种类别
$\{\mathsf{int},\mathsf{float},\mathsf{bool},\mathsf{event},\mathsf{view}\}$，
$\bar\tau$ 是有限类别序列；$\bar v:\mathsf{int}^*$ 简记 $\bar v$ 中各操作数均为 $\mathsf{int}$。
\begin{center}\small
\begin{tabularx}{\linewidth}{@{}lY@{}}
\toprule
定义 & 含义 \\
\midrule
$\Gamma:\mathsf{Var}\rightharpoonup\mathsf{Type}$
 & 当前作用域中已绑定变量的类别；$\Gamma,x:\tau$ 扩展一个新绑定。\\
$\mathsf{lit}_{\mathcal H}(c)=\tau$
 & 常量 $c$ 的唯一类别；不支持的常量无定义。\\
$\Delta_{\mathcal H}(g)=(\bar\tau,\tau)$
 & 标量算子 $g$ 的参数类别序列与结果类别；是部分签名表。\\
$\Sigma_{\mathcal H}(d,p,f)=\bar\tau$
 & 种类 $d$ 的管道 $p$ 可执行算子 $f$ 时的参数类别序列；无定义表示不支持该组合。\\
$L_{\mathcal H}\subseteq\mathsf{CoreKind}\times\mathsf{Pipe}\times\mathsf{Pipe}$
 & 允许的流内同步对 $(d,p,q)$，分别指定种类、源管道与目标管道。\\
$B_{\mathcal H}\subseteq\mathsf{Stream}\times\mathsf{Pipe}\times\mathsf{Stream}\times\mathsf{Pipe}$
 & 允许的流间同步边 $(s,p,t,q)$，从流 $s$ 的管道 $p$ 到流 $t$ 的管道 $q$。\\
$\mathsf{Barrier}_{\mathcal H}\subseteq\mathsf{CoreKind}\times\mathsf{Pipe}$
 & 各种类支持的屏障目标，包括配置允许的 $\mathsf{ALL}$。\\
\bottomrule
\end{tabularx}
\end{center}

在 \hyperlink{wf-sync}{\textsc{Sync}} 规则中，$a$ 仅取下表四种同步命令，
$v$ 必须是该命令自身的事件操作数。
$\mathsf{ok}_{\mathcal H,s}(a)$ 表示当前活动流中存在配置允许的同步端点：
\begin{center}\small
\begin{tabularx}{\linewidth}{@{}l l Y@{}}
\toprule
命令 $a$ & $\mathsf{ok}_{\mathcal H,s}(a)$ & 含义 \\
\midrule
$\kw{setflag}[p,q](v)$、$\kw{waitflag}[p,q](v)$
 & $(\mathsf{kind}(s),p,q)\in L_{\mathcal H}$ & 本流支持该同步对。\\
$\kw{syncblockset}[p,q](v)$
 & $\exists t\in\kappa.\ (s,p,t,q)\in B_{\mathcal H}$ & 至少有一条活动接收流。\\
$\kw{syncblockwait}[p,q](v)$
 & $\exists t\in\kappa.\ (t,p,s,q)\in B_{\mathcal H}$ & 至少有一条活动发送流。\\
\bottomrule
\end{tabularx}
\end{center}

\paragraph{配置归属。}
合法配置 $\mathcal H$ 须满足：
$\mathsf{SharedSpace}_{\mathcal H}\subseteq\mathsf{Space}$；
$(d,p,q)\in L_{\mathcal H}\Rightarrow p,q\in\mathsf{Pipe}_d$；
$(s,p,t,q)\in B_{\mathcal H}\Rightarrow s\ne t\land p\in\mathsf{Pipe}_{\mathsf{kind}(s)}\land q\in\mathsf{Pipe}_{\mathsf{kind}(t)}\cup\{\mathsf S\}$；
$(d,p)\in\mathsf{Barrier}_{\mathcal H}\Rightarrow p\in\mathsf{Pipe}_d\cup\{\mathsf{ALL}\}$。

\paragraph{跨核通知的配置。}
对本文的单轮广播配置，令
$T=\mathsf{Targets}_{\mathcal H}(\mathsf{aic},p,q)$。
$1{:}1$ 时 $T=\{\mathsf{aiv}_1\}$，$1{:}2$ 时
$T=\{\mathsf{aiv}_1,\mathsf{aiv}_2\}$。固定活动流与边后，令
\[
 D_h=\{t\in T\mid\beta_t(p,q,\varepsilon)=1\},
 \qquad h=(\mathsf{aic},p,q,\varepsilon).
\]
一次 \textsc{Block-Set} 完成使 $D_h=T$；每个接收者 $t$ 的等待完成仅移去 $t$。
该轮配对完成当且仅当置位后所有接收者均已消费，即 $D_h=\varnothing$。
第一个接收者可先继续执行；第二个接收者的消费决定该轮何时闭合。
同键下一轮置位完成须等上一轮闭合。程序结束时残留接收位表示基准的配对检查失败；
此检查独立于控制流能否返回。流内 flag 仍按 $(s,p,q,\varepsilon)$ 单独匹配。

此配置限制每个键至多有一轮在途。Atlas A2 的
\href{https://www.hiascend.com/doc_center/source/zh/canncommercial/80RC2/apiref/opdevgapi/atlasascendc_api_07_0268.html}{CrossCoreSetFlag 接口}
使用计数器，Cube 置位给两个 Vector 分别增加计数；反向汇聚需要两个 Vector 均置位。
反向汇聚及多轮重叠需要各自的配置与规则，本文只实例化上述正向广播。
\href{https://ascendnpu-ir.gitcode.com/en/developer_guide/features/auto_subtiling.html}{HIVM 自动子块划分}
受 \texttt{--enable-auto-bind-sub-block} 控制，可能回退至 $1{:}1$ 的工作划分。
构造基准时应分别记录编译目标、活动流以及该通知的实际参与者。

\litao{另待考虑：是否把 set/wait 的源端、目标端和 event 合并抽象为一个同步 identifier，
由配置解释其参与者与匹配关系，而非直接使用 $(p,q,\varepsilon)$。}

\begin{figure}[p]
\centering
\begin{tikzpicture}[
  x=1cm, y=0.84cm,
  font=\footnotesize,
  queue/.style={draw=#1, fill=white, rounded corners=2pt, minimum height=0.72cm, inner sep=1pt, align=center, font=\scriptsize\sffamily},
  scalarbox/.style={draw=scalarBorder, fill=scalarFill, rounded corners=2pt, align=center, font=\scriptsize},
  membox/.style 2 args={draw=#1, fill=#2, rounded corners=2pt, align=center, font=\scriptsize},
  dmaarr/.style={<->, >=Stealth, draw=gmBorder, thick},
  issuearr/.style={->, >=Stealth, draw=scalarBorder, thin},
  localarr/.style={->, >=Stealth, draw=syncLocal, densely dashed, semithick},
  crossarr/.style={->, >=Stealth, draw=syncCross, very thick},
  crossalt/.style={->, >=Stealth, draw=syncCross, densely dotted, thick}
]

% --- Top: Shared Global Memory ---
\draw[draw=gmBorder, fill=gmFill, rounded corners=3pt, thick] (0, 10.45) rectangle (15.4, 11.25);
\node[text=gmBorder!80!black] at (7.7, 10.85)
  {\textbf{共享全局内存}\; $\mathsf{SharedSpace}_{\mathcal H}=\{\mathsf{GM}\}$\quad（具体地址 $(\mathsf{GM}, z)$，全部活动流共享读写）};

% --- Outer AI Core container ---
\draw[draw=coreBorder, fill=coreBg, dashed, rounded corners=4pt, thick] (0, 0) rectangle (15.4, 9.65);
\node[fill=coreBg, inner sep=2.5pt, text=scalarBorder, font=\scriptsize\bfseries] at (3.43, 9.65)
  {单 AI Core（Ascend 910B $1{:}2$）};
\node[fill=coreBg, inner sep=2.5pt, text=scalarBorder, font=\scriptsize] at (11.82, 9.65)
  {活动流 $\kappa_{1:2}=\{\mathsf{aic},\,\mathsf{aiv}_1,\,\mathsf{aiv}_2\}$};

% --- Left Top: AIC Stream ---
\draw[draw=aicBorder, fill=aicFill, rounded corners=3pt, thick] (0.28, 4.30) rectangle (6.58, 9.20);
\node[anchor=west, text=aicBorder!90!black] at (0.45, 8.92)
  {\textbf{流 $s=\mathsf{aic}$}\;（$\mathsf{kind}(\mathsf{aic})=\mathsf{AIC}$）};

% AIC Scalar
\draw[scalarbox] (0.48, 8.02) rectangle (6.38, 8.64);
\node[font=\scriptsize] at (3.43, 8.33)
  {\textbf{Scalar 发射者}\; $(C_{\mathsf{aic}},\,\rho_{\mathsf{aic}})$\quad\color{scalarBorder}\scriptsize 控制流与 \kw{alloc}};

% AIC Queues Q_aic(p)
\node[queue=aicBorder, minimum width=1.24cm] (aic_mte2) at (1.10, 6.80) {\textbf{MTE2}\\{\tiny $Q_{\mathsf{aic}}(\mathsf{MTE2})$}};
\node[queue=aicBorder, minimum width=1.24cm] (aic_mte1) at (2.65, 6.80) {\textbf{MTE1}\\{\tiny $Q_{\mathsf{aic}}(\mathsf{MTE1})$}};
\node[queue=aicBorder, minimum width=1.24cm] (aic_m)    at (4.20, 6.80) {\textbf{M}\\{\tiny $Q_{\mathsf{aic}}(\mathsf{M})$}};
\node[queue=aicBorder, minimum width=1.24cm] (aic_fix)  at (5.75, 6.80) {\textbf{FIX}\\{\tiny $Q_{\mathsf{aic}}(\mathsf{FIX})$}};

% Scalar -> AIC queues
\foreach \q in {aic_mte2,aic_mte1,aic_m,aic_fix}{
  \draw[issuearr] (\q.north |- 0,8.02) -- (\q.north);
}

% Intra-AIC L_H arrows (separated anchors)
\draw[localarr] ([xshift=3pt]aic_mte2.south) -- ++(0,-0.25) -| ([xshift=-4pt]aic_mte1.south);
\draw[localarr] ([xshift=4pt]aic_mte1.south) -- ++(0,-0.25) -| ([xshift=-4pt]aic_m.south);
\draw[localarr] ([xshift=4pt]aic_m.south)    -- ++(0,-0.25) -| ([xshift=-3pt]aic_fix.south);
\node[fill=aicFill, inner sep=1pt, text=syncLocal, font=\scriptsize] at (3.43, 5.75)
  {$L_{\mathcal H}$：流内 \kw{setflag}/\kw{waitflag}（令牌 $\varphi_{\mathsf{aic}}$）};

% AIC Private On-Chip Memory
\draw[membox={aicBorder}{aicMem}] (0.48, 4.48) rectangle (6.38, 5.36);
\node[align=center, font=\scriptsize, text=aicBorder!85!black] at (3.43, 4.92)
  {\textbf{AIC 私有空间}：$\{\mathsf{L1},\,\mathsf{L0A},\,\mathsf{L0B},\,\mathsf{L0C}\}$\\
   地址 $(\mathsf{aic}, m, z)$};

% --- Right Top: AIV_1 Stream ---
\draw[draw=aivBorder, fill=aivFill, rounded corners=3pt, thick] (8.68, 4.95) rectangle (14.96, 9.20);
\node[anchor=west, text=aivBorder!90!black] at (8.85, 8.92)
  {\textbf{流 $s=\mathsf{aiv}_1$}\;（$\mathsf{kind}(\mathsf{aiv}_1)=\mathsf{AIV}$）};

% AIV_1 Scalar(S)
\draw[scalarbox] (8.88, 8.02) rectangle (14.76, 8.64);
\node[font=\scriptsize] at (11.82, 8.33)
  {\textbf{Scalar ($\mathsf{S}$) 发射者}\; $(C_{\mathsf{aiv}_1},\,\rho_{\mathsf{aiv}_1})$};

% AIV_1 Queues
\node[queue=aivBorder, minimum width=1.62cm] (aiv1_mte2) at (9.69, 6.80)  {\textbf{MTE2}\\{\tiny $Q_{\mathsf{aiv}_1}(\mathsf{MTE2})$}};
\node[queue=aivBorder, minimum width=1.62cm] (aiv1_v)    at (11.82, 6.80) {\textbf{V}\\{\tiny $Q_{\mathsf{aiv}_1}(\mathsf{V})$}};
\node[queue=aivBorder, minimum width=1.62cm] (aiv1_mte3) at (13.95, 6.80) {\textbf{MTE3}\\{\tiny $Q_{\mathsf{aiv}_1}(\mathsf{MTE3})$}};

\foreach \q in {aiv1_mte2,aiv1_v,aiv1_mte3}{
  \draw[issuearr] (\q.north |- 0,8.02) -- (\q.north);
}

% Barrier_H marker on AIV_1 (V queue + ALL line)
\draw[draw=syncBar, very thick] (12.34, 6.47) -- (12.34, 7.13);
\draw[draw=syncBar, very thick] (12.44, 6.47) -- (12.44, 7.13);
\draw[draw=syncBar, densely dotted, thick] (8.88, 6.26) -- (14.76, 6.26);
\node[fill=aivFill, inner sep=1pt, text=syncBar, font=\tiny] at (12.88, 6.26) {\kw{bar}[$\mathsf{ALL}$]};

% Intra-AIV_1 L_H arrows (separated anchors)
\draw[localarr] ([xshift=4pt]aiv1_mte2.south) -- ++(0,-0.32) -| ([xshift=-5pt]aiv1_v.south);
\draw[localarr] ([xshift=5pt]aiv1_v.south)    -- ++(0,-0.32) -| ([xshift=-4pt]aiv1_mte3.south);
\node[fill=aivFill, inner sep=1pt, text=syncLocal, font=\scriptsize] at (11.82, 5.86)
  {$L_{\mathcal H}\ (\varphi_{\mathsf{aiv}_1})$\; /\; \color{syncBar}$\mathsf{Barrier}_{\mathcal H}$};

% AIV_1 Private UB
\draw[membox={aivBorder}{aivMem}] (8.88, 5.10) rectangle (14.76, 5.64);
\node[font=\scriptsize, text=aivBorder!85!black] at (11.82, 5.37)
  {\textbf{$\mathsf{aiv}_1$ 私有 $\mathsf{UB}$}：地址 $(\mathsf{aiv}_1, \mathsf{UB}, z)$};

% --- Right Bottom: AIV_2 Stream ---
\draw[draw=aivBorder, fill=aivFill, rounded corners=3pt, thick] (8.68, 0.28) rectangle (14.96, 4.53);
\node[anchor=west, text=aivBorder!90!black] at (8.85, 4.25)
  {\textbf{流 $s=\mathsf{aiv}_2$}\;（$\mathsf{kind}(\mathsf{aiv}_2)=\mathsf{AIV}$）};

% AIV_2 Scalar(S)
\draw[scalarbox] (8.88, 3.35) rectangle (14.76, 3.97);
\node[font=\scriptsize] at (11.82, 3.66)
  {\textbf{Scalar ($\mathsf{S}$) 发射者}\; $(C_{\mathsf{aiv}_2},\,\rho_{\mathsf{aiv}_2})$};

% AIV_2 Queues
\node[queue=aivBorder, minimum width=1.62cm] (aiv2_mte2) at (9.69, 2.21)  {\textbf{MTE2}\\{\tiny $Q_{\mathsf{aiv}_2}(\mathsf{MTE2})$}};
\node[queue=aivBorder, minimum width=1.62cm] (aiv2_v)    at (11.82, 2.21) {\textbf{V}\\{\tiny $Q_{\mathsf{aiv}_2}(\mathsf{V})$}};
\node[queue=aivBorder, minimum width=1.62cm] (aiv2_mte3) at (13.95, 2.21) {\textbf{MTE3}\\{\tiny $Q_{\mathsf{aiv}_2}(\mathsf{MTE3})$}};

\foreach \q in {aiv2_mte2,aiv2_v,aiv2_mte3}{
  \draw[issuearr] (\q.north |- 0,3.35) -- (\q.north);
}

\draw[draw=syncBar, very thick] (12.34, 1.88) -- (12.34, 2.54);
\draw[draw=syncBar, very thick] (12.44, 1.88) -- (12.44, 2.54);
\draw[localarr] ([xshift=4pt]aiv2_mte2.south) -- ++(0,-0.28) -| ([xshift=-5pt]aiv2_v.south);
\draw[localarr] ([xshift=5pt]aiv2_v.south)    -- ++(0,-0.28) -| ([xshift=-4pt]aiv2_mte3.south);
\node[fill=aivFill, inner sep=1pt, text=syncLocal, font=\scriptsize] at (11.82, 1.28)
  {$L_{\mathcal H}\ (\varphi_{\mathsf{aiv}_2})$\; /\; \color{syncBar}$\mathsf{Barrier}_{\mathcal H}$};

% AIV_2 Private UB
\draw[membox={aivBorder}{aivMem}] (8.88, 0.45) rectangle (14.76, 0.99);
\node[font=\scriptsize, text=aivBorder!85!black] at (11.82, 0.72)
  {\textbf{$\mathsf{aiv}_2$ 私有 $\mathsf{UB}$}：$(\mathsf{aiv}_2, \mathsf{UB}, z)$（与 $\mathsf{aiv}_1$ 隔离）};

% --- GM <-> Streams DMA links ---
\draw[->, >=Stealth, draw=gmBorder, thick] (1.10, 10.45) -- (1.10, 9.20);
\node[right, font=\tiny, text=gmBorder!85!black] at (1.12, 10.08) {MTE2};
\draw[->, >=Stealth, draw=gmBorder, thick] (5.75, 9.20) -- (5.75, 10.45);
\node[left, font=\tiny, text=gmBorder!85!black] at (5.73, 10.08) {FIX};
\draw[->, >=Stealth, draw=gmBorder, thick] (9.69, 10.45) -- (9.69, 9.20);
\node[right, font=\tiny, text=gmBorder!85!black] at (9.71, 10.08) {MTE2};
\draw[->, >=Stealth, draw=gmBorder, thick] (13.95, 9.20) -- (13.95, 10.45);
\node[right, font=\tiny, text=gmBorder!85!black] at (13.97, 10.08) {MTE3};
% Right-margin bus to AIV_2
\draw[dmaarr] (15.18, 10.45) |- (14.96, 2.21);
\node[rotate=90, font=\tiny, text=gmBorder!85!black, fill=coreBg, inner sep=1pt] at (15.18, 4.74) {GM 访存};

% --- Cross-Core 1:2 Broadcast B_H ---
\coordinate (bcast) at (7.63, 6.80);
\draw[draw=syncCross, very thick] (aic_fix.east) -- (bcast);
\fill[syncCross] (bcast) circle (2.2pt);
\draw[crossarr] (bcast) |- (8.88, 8.33);
\draw[crossarr] (bcast) |- (8.88, 3.66);
\draw[crossalt] (bcast) -- (aiv1_mte2.west);

\node[fill=coreBg, draw=syncCross!50, rounded corners=2pt, inner sep=2.5pt, align=center, font=\scriptsize, text=syncCross] at (7.63, 5.20)
  {\textbf{$B_{\mathcal H}$ 跨核 1:2}\\$\beta_{\mathsf{aiv}_1},\,\beta_{\mathsf{aiv}_2}$};
\node[fill=coreBg, inner sep=1pt, font=\tiny, text=syncCross] at (8.20, 8.53) {$q=\mathsf{S}$};
\node[fill=coreBg, inner sep=1pt, font=\tiny, text=syncCross] at (8.20, 3.86) {$q=\mathsf{S}$};
\node[fill=coreBg, inner sep=1pt, font=\tiny, text=syncCross] at (8.25, 7.02) {$q\ne\mathsf{S}$};

% --- Bottom Left: Legend / Parameter Mapping ---
\draw[draw=coreBorder, fill=white, rounded corners=3pt] (0.28, 0.28) rectangle (6.58, 3.95);
\node[anchor=west, font=\scriptsize\bfseries, text=scalarBorder] at (0.45, 3.66)
  {图例：硬件参数类型与形式化抽象对应};
\draw[draw=coreBorder!60] (0.42, 3.44) -- (6.44, 3.44);

% Legend rows
\filldraw[fill=gmFill, draw=gmBorder] (0.48, 3.04) rectangle (0.90, 3.26);
\node[anchor=west, font=\scriptsize] at (1.00, 3.15) {共享空间 $\mathsf{SharedSpace}_{\mathcal H}$：$\mathsf{GM}$，$(m,z)$};

\filldraw[fill=aicFill, draw=aicBorder] (0.48, 2.60) rectangle (0.67, 2.82);
\filldraw[fill=aivFill, draw=aivBorder] (0.71, 2.60) rectangle (0.90, 2.82);
\node[anchor=west, font=\scriptsize] at (1.00, 2.71) {核类 $\mathsf{CoreKind}$ / 流 $s\in\kappa$：私有 $(s,m,z)$};

\filldraw[fill=scalarFill, draw=scalarBorder] (0.48, 2.16) rectangle (0.90, 2.38);
\node[anchor=west, font=\scriptsize] at (1.00, 2.27) {标量发射者 $(C_s,\rho_s)$：控制流与 \kw{alloc} ($\mathsf{S}$)};

\filldraw[fill=white, draw=scalarBorder] (0.48, 1.72) rectangle (0.90, 1.94);
\node[anchor=west, font=\scriptsize] at (1.00, 1.83) {管道队列 $Q_s(p)$：$p\in\mathsf{Pipe}_d$，异步乱序完成};

\draw[localarr] (0.46, 1.39) -- (0.92, 1.39);
\node[anchor=west, font=\scriptsize] at (1.00, 1.39) {流内同步 $L_{\mathcal H}$：\kw{setflag}/\kw{waitflag}（令牌 $\varphi_s$）};

\draw[draw=syncBar, very thick] (0.64, 0.84) -- (0.64, 1.08);
\draw[draw=syncBar, very thick] (0.74, 0.84) -- (0.74, 1.08);
\node[anchor=west, font=\scriptsize] at (1.00, 0.96) {流内屏障 $\mathsf{Barrier}_{\mathcal H}$：\kw{pipebarrier}$[p]$ 或 $\mathsf{ALL}$};

\draw[crossarr] (0.46, 0.54) -- (0.92, 0.54);
\node[anchor=west, font=\scriptsize] at (1.00, 0.54) {跨核同步 $B_{\mathcal H}$：\kw{syncblock} 1:2 广播（$\beta_t$）};

\end{tikzpicture}
\caption{静态硬件配置与形式化抽象的对应（以 Ascend 910B $1{:}2$ 配置为例）：共享内存 $\mathsf{SharedSpace}_{\mathcal H}$、核类与私有地址空间 $\mathsf{Loc}_{s,m}$、数据管道 $\mathsf{Pipe}_d$ 及三类同步关系 $L_{\mathcal H},B_{\mathcal H},\mathsf{Barrier}_{\mathcal H}$。}
\label{fig:arch-mapping}

\vspace{0.45cm}

\begin{tikzpicture}[
  x=1cm, y=0.88cm,
  font=\footnotesize,
  panel/.style 2 args={draw=#1, fill=#2, rounded corners=3pt, thick},
  openitem/.style={draw=aivBorder, fill=aivFill, rounded corners=2pt, minimum height=0.64cm, inner sep=2pt, align=center, font=\scriptsize},
  blockeditem/.style={draw=coreBorder, fill=scalarFill, rounded corners=2pt, minimum height=0.64cm, inner sep=2pt, align=center, font=\scriptsize, text=scalarBorder},
  syncitem/.style={draw=syncLocal, fill=aicFill, rounded corners=2pt, minimum height=0.64cm, inner sep=2pt, align=center, font=\scriptsize, text=syncLocal},
  baritem/.style={draw=syncBar, fill=red!6, rounded corners=2pt, minimum height=0.64cm, inner sep=2pt, align=center, font=\scriptsize, text=syncBar},
  crossitem/.style={draw=syncCross, fill=syncCross!10, rounded corners=2pt, minimum height=0.64cm, inner sep=2pt, align=center, font=\scriptsize, text=syncCross},
  steparr/.style={->, >=Stealth, semithick}
]

% --- Left Column: Instruction Stream C_s ---
\draw[panel={scalarBorder}{coreBg}] (0, 0) rectangle (4.35, 8.75);
\node[anchor=west, font=\scriptsize\bfseries, text=scalarBorder] at (0.18, 8.42)
  {单条指令流 $s$：控制端 $(C_s,\rho_s)$};
\draw[draw=coreBorder!70] (0.15, 8.14) -- (4.20, 8.14);

\node[draw=scalarBorder, fill=white, rounded corners=2pt, text width=3.75cm, align=left, font=\scriptsize, inner sep=3.5pt] (cmd1) at (2.175, 7.12)
  {\textbf{1. 标量 / 分配 / 控制}\\
   \kw{let}\ $x=\kw{scalar}[g](\bar v)$\\
   \kw{let}\ $x=\kw{alloc}[m](\bar v)$\\
   \kw{if} / \kw{for}（立即更新 $\rho_s,\mu$）};

\node[draw=aicBorder, fill=aicFill!70, rounded corners=2pt, text width=3.75cm, align=left, font=\scriptsize, inner sep=3.5pt] (cmd2) at (2.175, 4.72)
  {\textbf{2. 异步发射（非阻塞）}\\
   \kw{issue}$[p](f,\bar v)$\\
   \kw{setflag} / \kw{waitflag} / \kw{pipebarrier}\\
   \kw{syncblockset} / \kw{bwait}($q\!\ne\!\mathsf S$)\\
   \color{aicBorder!90!black}$\Rightarrow$ 追加至队尾，控制立即继续};

\node[draw=syncCross, fill=syncCross!8, rounded corners=2pt, text width=3.75cm, align=left, font=\scriptsize, inner sep=3.5pt] (cmd3) at (2.175, 2.28)
  {\textbf{3. 标量阻塞等待}\\
   \kw{syncblockwait}$[p,\mathsf S](v)$\\
   \color{syncCross}$\Rightarrow$ \textsc{Scalar-Wait}：\\
   阻塞本流控制直至 $\beta_s=1$};

\draw[steparr, draw=scalarBorder] (cmd1.south) -- (cmd2.north);
\draw[steparr, draw=scalarBorder] (cmd2.south) -- (cmd3.north);

\node[text width=3.95cm, align=center, font=\tiny, text=scalarBorder] at (2.175, 0.65)
  {全局小步关系 $\longrightarrow_{\mathcal H}$ 交错推进各流 $s\in\kappa$ 的控制步与队列完成步};

% --- Top-Right Panel: Async Pipe Queues Q_s(p) ---
\draw[panel={aicBorder!75}{white}] (4.80, 2.85) rectangle (15.40, 8.75);
\node[anchor=west, font=\scriptsize\bfseries, text=aicBorder!90!black] at (4.98, 8.42)
  {异步管道队列 $Q_s(p)$（队首在左 $\leftarrow$ 移出完成；队尾在右 $\leftarrow$ 发射入队）};
\draw[draw=aicBorder!35] (4.95, 8.14) -- (15.25, 8.14);

% Queue 1: MTE2
\node[anchor=west, font=\scriptsize\sffamily\bfseries, text=aivBorder!90!black] at (4.98, 7.82) {$Q_s(\mathsf{MTE2})$};
\draw[draw=coreBorder!60, fill=coreBg, rounded corners=2pt] (4.98, 6.80) rectangle (14.65, 7.62);

\node[openitem, minimum width=2.45cm] (q1_op1) at (6.32, 7.21) {\kw{op}$_1(R_1,W_1)$\\[-2pt]{\tiny 已打开 $\mathsf{open}$}};
\node[openitem, minimum width=2.45cm] (q1_op2) at (8.98, 7.21) {\kw{op}$_2(R_2,W_2)$\\[-2pt]{\tiny 已打开（可乱序）}};
\node[syncitem, minimum width=2.75cm] (q1_set) at (11.78, 7.21) {\kw{set}$(\mathsf V,\varepsilon)$\\[-2pt]{\tiny 须到队首且 $\varphi_s{=}0$}};
\node[font=\tiny, text=aicBorder!90!black] at (13.95, 7.21) {$\leftarrow$ 入队};

% Queue 2: V
\node[anchor=west, font=\scriptsize\sffamily\bfseries, text=aivBorder!90!black] at (4.98, 6.26) {$Q_s(\mathsf{V})$};
\draw[draw=coreBorder!60, fill=coreBg, rounded corners=2pt] (4.98, 5.24) rectangle (14.65, 6.06);

\node[openitem, minimum width=1.60cm]    (q2_opa)  at (5.90, 5.65) {\kw{op}$_a$\\[-2pt]{\tiny 已打开}};
\node[syncitem, minimum width=2.15cm]    (q2_wait) at (7.92, 5.65) {\kw{wait}$(\mathsf{MTE2},\varepsilon)$\\[-2pt]{\tiny $\mathsf{open}(U)\land\varphi_s{=}1$}};
\node[blockeditem, minimum width=1.65cm] (q2_opb)  at (9.96, 5.65) {\kw{op}$_b$\\[-2pt]{\tiny 被 \kw{wait} 遮挡}};
\node[baritem, minimum width=1.22cm]     (q2_bar)  at (11.54, 5.65) {\kw{bar}\\[-2pt]{\tiny 须到队首}};
\node[blockeditem, minimum width=1.55cm] (q2_opc)  at (13.06, 5.65) {\kw{op}$_c$\\[-2pt]{\tiny 被 \kw{bar} 遮挡}};
\node[font=\tiny, text=aicBorder!90!black] at (14.24, 5.65) {$\leftarrow$};

% Queue 3: FIX
\node[anchor=west, font=\scriptsize\sffamily\bfseries, text=aicBorder!90!black] at (4.98, 4.70) {$Q_{\mathsf{aic}}(\mathsf{FIX})$};
\draw[draw=coreBorder!60, fill=coreBg, rounded corners=2pt] (4.98, 3.68) rectangle (14.65, 4.50);

\node[openitem, minimum width=2.45cm]  (q3_op)   at (6.32, 4.09) {\kw{op}$(\kw{Fixpipe})$\\[-2pt]{\tiny 已打开 $\mathsf{open}$}};
\node[crossitem, minimum width=3.85cm] (q3_bset) at (9.75, 4.09) {\kw{bset}$(\mathsf S,\varepsilon')$\\[-2pt]{\tiny 到队首且 $\forall t\in T.\,\beta_t{=}0$}};
\node[font=\tiny, text=aicBorder!90!black] at (13.95, 4.09) {$\leftarrow$ 入队};

% Rule highlight bar inside Queue panel (compact left block, leaves right clear for bset arrow)
\node[anchor=west, text width=7.55cm, align=left, font=\tiny, text=scalarBorder] at (4.98, 3.22)
  {\textcolor{aivBorder}{\textbf{绿底（已打开）}}：前缀不含 \kw{wait}/\kw{bwait}/\kw{bar}，可乱序 \textsc{Op-Done}；\\
   \textcolor{scalarBorder}{\textbf{灰底（未打开）}}：被前方 \kw{wait}/\kw{bar} 遮挡，不可提前完成。};

% --- Clean Enqueue Bus: cmd2 -> right (tail) of queues along top-right spine ---
\draw[draw=aicBorder, semithick] (cmd2.east) -- (4.58, 4.72) |- (4.58, 8.95) -- (15.02, 8.95) -- (15.02, 4.09);
\draw[steparr, draw=aicBorder] (15.02, 7.21) -- (14.65, 7.21);
\draw[steparr, draw=aicBorder] (15.02, 5.65) -- (14.65, 5.65);
\draw[steparr, draw=aicBorder] (15.02, 4.09) -- (14.65, 4.09);
\node[fill=white, inner sep=1.5pt, font=\tiny, text=aicBorder!90!black] at (9.90, 8.95)
  {\textsc{Issue} / \textsc{Sync-Issue} / \textsc{Barrier-Issue}（非阻塞追加至目标队列队尾）};

% --- Bottom Right: Shared State (Memory + Tokens) ---
% Shared Memory Box
\draw[panel={gmBorder}{gmFill}] (4.80, 0) rectangle (10.10, 2.15);
\node[anchor=west, font=\scriptsize\bfseries, text=gmBorder!85!black] at (4.96, 1.86)
  {共享内存 $\mu:\mathsf{Loc}\rightharpoonup\mathsf{Val}\uplus\{\mathsf{Uninit}\}$};
\draw[draw=gmBorder!35] (4.94, 1.64) -- (9.96, 1.64);
\node[anchor=north west, text width=4.95cm, align=left, font=\tiny, text=gmBorder!65!black] at (4.96, 1.56)
  {\textbullet\;\textsc{Alloc}：新鲜集 $A\cap\mathrm{dom}(\mu)=\varnothing$，置 $\mathsf{Uninit}$\\[2pt]
   \textbullet\;\textsc{Op-Done}：检查 $R\cup W\subseteq\mathrm{dom}(\mu)\cap\mathsf{Loc}_s$，更新 $\mu\oplus\delta$\\[2pt]
   \textbullet\;\textcolor{syncBar}{\textbf{Data Race}}：两已打开 \kw{op} 的 $(R,W)$ 足迹冲突};

% Token Pool Box
\draw[panel={syncCross!80}{syncCross!5}] (10.42, 0) rectangle (15.40, 2.15);
\node[anchor=west, font=\scriptsize\bfseries, text=syncCross] at (10.58, 1.86)
  {同步令牌状态（初始全 $0$）};
\draw[draw=syncCross!30] (10.56, 1.64) -- (15.26, 1.64);
\node[anchor=north west, text width=4.65cm, align=left, font=\tiny, text=scalarBorder] at (10.58, 1.56)
  {\textbullet\;\textcolor{syncLocal}{\textbf{流内 $\varphi_s$}}：\textsc{Set-Done} 置 $1$，\textsc{Wait-Done} 清 $0$\\[3pt]
   \textbullet\;\textcolor{syncCross}{\textbf{跨核 $\beta_t$}}：\textsc{Block-Set} 置全 $t\in T$ 为 $1$；\\
   \quad 各流 \textsc{Scalar}/\textsc{Block-Wait} 独立清 $0$};

% --- Semantic flow links (clean corridor at y = 2.50, zero box/text crossing) ---
% Op-Done -> Memory
\draw[steparr, draw=gmBorder, densely dashed] (5.55, 2.85) -- (5.55, 2.15);
\node[right, font=\tiny, text=gmBorder!85!black] at (5.58, 2.33) {\textsc{Op-Done}};
% Intra-stream set -> wait (inside Queue panel, routed between Q1 and Q2)
\draw[steparr, draw=syncLocal, densely dashed] (q1_set.south) -- ++(0,-0.32) -| (q2_wait.north);
\node[fill=white, inner sep=1pt, font=\tiny, text=syncLocal] at (10.45, 6.46) {$\varphi_s=1$};
% Cross-core bset -> Token Pool (routed cleanly down right-side open space)
\draw[steparr, draw=syncCross, thick] (q3_bset.east) -- (12.90, 4.09) -- (12.90, 2.15);
\node[fill=white, inner sep=1pt, font=\tiny, text=syncCross] at (13.95, 2.52) {\textsc{Block-Set} ($\beta_t\gets 1$)};
% Token Pool -> Scalar-Wait (routed along clean upper corridor y = 2.60 to cmd3.east!)
\draw[steparr, draw=syncCross, thick] (10.42, 1.92) -- (10.26, 1.92) |- (4.35, 2.60);
\node[fill=white, inner sep=1pt, font=\tiny, text=syncCross] at (8.35, 2.60) {$\beta_s=1$ 唤醒};

\end{tikzpicture}
\caption{小步操作语义的动态执行模型（与具体芯片拓扑解耦）：标量控制端将异步操作与同步原语非阻塞地追加到各管道队列 $Q_s(p)$ 队尾（仅 \textsc{Scalar-Wait} 阻塞控制端）；队列中满足 $\mathsf{open}(U)$ 的普通操作可乱序完成并原子读写 $\mu$，\kw{set}/\kw{bset}/\kw{bar} 须位于队首方可完成并更新令牌 $\varphi_s,\beta_t$。}
\label{fig:opsem-exec}
\end{figure}

\section{状态与参数}

固定配置 $\mathcal H$ 与活动流集合 $\kappa$。运行时状态及动态解释如下，
其中 $\mathsf{Val}$ 为值域，$\mathsf{Cell}=\mathsf{Val}\uplus\{\mathsf{Uninit}\}$。

\begin{center}\small
\begin{tabularx}{\linewidth}{@{}lY@{}}
\toprule
状态分量 & 含义 \\
\midrule
$\mu:\mathsf{Loc}\rightharpoonup\mathsf{Cell}$
 & 共享内存；已分配单元保存值或未初始化标记。\\
$\rho_s:\mathsf{Var}\rightharpoonup\mathsf{Val}$
 & 指令流 $s$ 的局部变量环境；操作数求值时常量映到自身。\\
$Q_s:\mathsf{Pipe}_{\mathsf{kind}(s)}\to\mathsf{Item}^{*}$
 & 指令流 $s$ 各数据管道的有限队列；队列项见下式。\\
$\varphi_s(p,q,\varepsilon)\in\{0,1\}$
 & 流内同步令牌，记录源管道 $p$ 到目标管道 $q$ 的通知。\\
$\beta_t(p,q,\varepsilon)\in\{0,1\}$
 & 块同步令牌，按接收流 $t$ 分开记录；与流内令牌独立。\\
\midrule
动态参数与记号 & 含义 \\
\midrule
$\mathsf{eval}_{\mathcal H}(s,g,\bar\nu)=\nu$
 & 标量与视图构造的部分求值函数；成功结果符合签名 $\Delta_{\mathcal H}$。\\
$\mathsf{alloc}_{\mathcal H}(s,m,\bar\nu,A,\nu)$
 & 在空间 $m$ 中按维度参数 $\bar\nu$ 分配有限具体地址集 $A\subseteq\mathsf{Loc}_{s,m}$ 并给出对应视图 $\nu$ 的关系；参数非法或超容量时为空。\\
$\mathsf{foot}_{\mathcal H}(s,p,f,\bar\nu)=(R,W)$
 & 将实际 view 解释为有限读区 $R$ 与写区 $W$，包含寻址、布局及别名信息。\\
$F_{\mathcal H}(s,p,f,\bar\nu,\mu|_R,\delta)$
 & 根据操作数及完成时的读区内容产生写入更新 $\delta$ 的关系；允许多个结果。\\
$\delta:\mathsf{Loc}\rightharpoonup\mathsf{Cell}$，$\mathrm{dom}(\delta)\subseteq W$
 & 有限内存更新；$\mu\oplus\delta$ 仅覆盖 $\delta$ 定义的单元。\\
$\mathsf{Targets}_{\mathcal H}(s,p,q)$
 & 由同步边 $B_{\mathcal H}$ 导出的活动接收流集合，定义见下式。\\
\bottomrule
\end{tabularx}
\end{center}

队列项 $\iota\in\mathsf{Item}$ 与接收流集合定义为
\[
\begin{aligned}
\iota ::={}
  &\ \kw{op}(f,\bar\nu,R,W)
   \mid \kw{set}(q,\varepsilon)
   \mid \kw{wait}(p,\varepsilon)\\
  &\mid \kw{bar}
   \mid \kw{bset}(q,\varepsilon)
   \mid \kw{bwait}(p,\varepsilon),\\[0.4em]
\mathsf{Targets}_{\mathcal H}(s,p,q)
  & =\{t\in\kappa\mid(s,p,t,q)\in B_{\mathcal H}\}.
\end{aligned}
\]
\begin{center}\small
\begin{tabularx}{\linewidth}{@{}l l Y c@{}}
\toprule
队列项 $\iota$ & 来源与所在队列 & 完成条件与效果 & 遮挡后续 \\
\midrule
$\kw{op}(f,\bar\nu,R,W)$ & $\kw{issue}[p]\to Q_s(p)$ & $\mathsf{open}(U)$，原子读 $\mu|_R$ 并写 $\mu\oplus\delta$ & 否 \\
$\kw{set}(q,\varepsilon)$ & $\kw{setflag}[p,q]\to Q_s(p)$ & 队首且 $\varphi_s(p,q,\varepsilon)=0$，置为 $1$ & 否 \\
$\kw{wait}(p,\varepsilon)$ & $\kw{waitflag}[p,q]\to Q_s(q)$ & $\mathsf{open}(U)$ 且 $\varphi_s(p,q,\varepsilon)=1$，清为 $0$ & 是 \\
$\kw{bar}$ & $\kw{pipebarrier}[p]\to Q_s(p)$（$\mathsf{ALL}$ 入全队列） & 队首直接移出（排空前序操作） & 是 \\
$\kw{bset}(q,\varepsilon)$ & $\kw{syncblockset}[p,q]\to Q_s(p)$ & 队首且 $\forall t\in T.\,\beta_t=0$，全置为 $1$ & 否 \\
$\kw{bwait}(p,\varepsilon)$ & $\kw{syncblockwait}[p,q]$（$q\ne\mathsf S$）$\to Q_s(q)$ & $\mathsf{open}(U)$ 且本流 $\beta_s=1$，清为 $0$ & 是 \\
\bottomrule
\end{tabularx}
\end{center}
$\kw{syncblockwait}[p,\mathsf S]$ 不入队，直接由 \textsc{Scalar-Wait} 阻塞控制端。
队列保留实际操作数 $\bar\nu$，以保留足迹集合无法表达的复制对应、步幅和标量参数。
初始状态 $\sigma_0$ 中，各条流的队列 $Q_s$ 均为空，所有流内令牌 $\varphi_s$ 与块同步令牌 $\beta_t$ 初值均为 $0$，初始内存 $\mu_0$ 包含核启动前由外部传入的已分配输入输出存储区；块置位通知各接收流，等待只消费本流的令牌。

执行配置与小步关系为
\[
\begin{gathered}
\langle P,\sigma\rangle,\qquad
P=\mathop{\Big\|}_{s\in\kappa}s\{C_s\},\qquad
\sigma=\langle\mu,\beta,(\rho_s,Q_s,\varphi_s)_{s\in\kappa}\rangle,\\
\langle P,\sigma\rangle\longrightarrow_{\mathcal H}\langle P',\sigma'\rangle.
\end{gathered}
\]
$C_s$ 直接表示当前待执行的命令，$\sigma$ 不另存剩余程序。
控制步改写程序项；管道操作完成时可以只更新状态。$\mathcal H$ 不随执行改变。

队列中位于等待项或屏障项之后的项尚未打开。$\kw{op}$ 在自身已经打开时可以完成，不要求位于队首。$\kw{set}$、$\kw{bset}$ 与 $\kw{bar}$ 只有在前面的项都已离开队列时才能完成。

\clearpage
\section{小步规则}

以 $\langle C,\sigma\rangle\xrightarrow{s}_{\mathcal H}\langle C',\sigma'\rangle$
表示流 $s$ 的一步；$\sigma$ 仍是全局状态。记 $\rho=\rho_s$，
$\sigma[X\mapsto a]$ 仅更新分量 $X$，未列出的分量保持不变。
以下规则归纳定义小步关系（取满足规则的最小关系）；全局执行交错各流的步骤：
\[
\infer{\langle C_s,\sigma\rangle\xrightarrow{s}_{\mathcal H}\langle C_s',\sigma'\rangle}
 {\langle s\{C_s\}\parallel P_{\rm rest},\sigma\rangle
 \longrightarrow_{\mathcal H}
 \langle s\{C_s'\}\parallel P_{\rm rest},\sigma'\rangle}
 \;\textsc{Step}
\]

\begin{figure}[!htbp]
\centering\small
\[
\begin{gathered}
\infer{\rho(\bar v)=\bar\nu\quad\mathsf{eval}_{\mathcal H}(s,g,\bar\nu)=\nu}
 {\langle\kw{let}\ x=\kw{scalar}[g](\bar v)\ \kw{in}\ C,\sigma\rangle
 \xrightarrow{s}_{\mathcal H}\langle C,\sigma[\rho_s\mapsto\rho[x\mapsto\nu]]\rangle}
 \;\textsc{Scalar}\\[1em]
\infer{\rho(\bar v)=\bar\nu\quad\mathsf{alloc}_{\mathcal H}(s,m,\bar\nu,A,\nu)\quad A\cap\mathrm{dom}(\mu)=\varnothing}
 {\langle\kw{let}\ x=\kw{alloc}[m](\bar v)\ \kw{in}\ C,\sigma\rangle
 \xrightarrow{s}_{\mathcal H}\langle C,\sigma[\mu\mapsto\mu\uplus(A\mapsto\mathsf{Uninit}),\ \rho_s\mapsto\rho[x\mapsto\nu]]\rangle}
 \;\textsc{Alloc}\\[1em]
\infer{\rho(v)=b\quad b\in\{\kw{true},\kw{false}\}}
 {\langle\kw{if}\ v\ \kw{then}\ \{C_{\kw{true}}\}\ \kw{else}\ \{C_{\kw{false}}\};\ C_3,\sigma\rangle
 \xrightarrow{s}_{\mathcal H}\langle C'_b\triangleright C_3,\sigma\rangle}
 \;\textsc{Branch}\\[1em]
\infer{\rho(v_i)=k_i\ (i=1,2,3)\quad\mathsf{more}(k_1,k_2,k_3)}
 {\langle\mathsf{Loop}(x,\bar v,B,C),\sigma\rangle
 \xrightarrow{s}_{\mathcal H}\langle B'[k_1/x]\triangleright
 \mathsf{Loop}(x,(k_1+k_3,k_2,k_3),B,C),\sigma\rangle}
 \;\textsc{For-Step}\\[1em]
\infer{\rho(v_i)=k_i\ (i=1,2,3)\quad k_3\ne0\quad\neg\mathsf{more}(k_1,k_2,k_3)}
 {\langle\mathsf{Loop}(x,\bar v,B,C),\sigma\rangle
 \xrightarrow{s}_{\mathcal H}\langle C,\sigma\rangle}
 \;\textsc{For-Done}
\end{gathered}
\]
\caption{标量控制、内存分配与循环}
\begin{itemize}
\item $\mu\uplus(A\mapsto\mathsf{Uninit})$ 将新鲜地址集 $A$ 中各单元初始化为 $\mathsf{Uninit}$ 并并入 $\mu$；$A\cap\mathrm{dom}(\mu)=\varnothing$ 保证每次 \kw{alloc} 互不别名。
\item $\mathsf{Loop}(x,\bar v,B,C)$ 缩写
$\kw{for}\ x=v_1\ \kw{to}\ v_2\ \kw{step}\ v_3\ \kw{do}\ \{B\};C$；
$\mathsf{more}(a,b,h)\iff(h>0\land a<b)\lor(h<0\land a>b)$。
\item \textsc{Branch} 与 \textsc{For-Step} 拼接续延前均对内层体中局部绑定变量作新鲜 $\alpha$-重命名（记作 $C'_b,B'$），\textsc{For-Step} 另以无捕获代换绑定本轮循环变量 $k_1$。
$B\triangleright D$ 将本层的 \kw{end} 替换为 $D$；保留 \kw{return}，递归处理后续命令，遇到嵌套 \kw{if} 或 \kw{for} 只替换其续延，不进入其内层分支或循环体。
\item \kw{return} 与 \kw{end} 没有控制步骤，但已入队操作仍可完成。零步长、未定义标量求值或分配失败没有对应控制步骤。
\end{itemize}
\end{figure}

\clearpage
\begin{figure}[!htbp]
\centering\small
\[
\begin{gathered}
\infer{\rho(\bar v)=\bar\nu\quad\mathsf{foot}_{\mathcal H}(s,p,f,\bar\nu)=(R,W)}
 {\langle\kw{issue}[p](f,\bar v);C,\sigma\rangle\xrightarrow{s}_{\mathcal H}
 \langle C,\sigma[Q_s(p)\mapsto Q_s(p)\cdot\kw{op}(f,\bar\nu,R,W)]\rangle}
 \;\textsc{Issue}\\[1em]
\infer{\rho(v)=\varepsilon\quad(a,r,\iota)\in\mathsf{Enqueue}(p,q,v,\varepsilon)}
 {\langle a;C,\sigma\rangle\xrightarrow{s}_{\mathcal H}
 \langle C,\sigma[Q_s(r)\mapsto Q_s(r)\cdot\iota]\rangle}
 \;\textsc{Sync-Issue}\\[1em]
\infer{D=\begin{cases}\mathsf{Pipe}_{\mathsf{kind}(s)},&p=\mathsf{ALL},\\\{p\},&p\in\mathsf{Pipe}_{\mathsf{kind}(s)}\end{cases}}
 {\langle\kw{pipebarrier}[p];C,\sigma\rangle\xrightarrow{s}_{\mathcal H}
 \langle C,\sigma[Q_s(r)\mapsto Q_s(r)\cdot\kw{bar}]_{r\in D}\rangle}
 \;\textsc{Barrier-Issue}\\[1em]
\infer{\rho(v)=\varepsilon\quad\beta_s(p,\mathsf S,\varepsilon)=1}
 {\langle\kw{syncblockwait}[p,\mathsf S](v);C,\sigma\rangle\xrightarrow{s}_{\mathcal H}
 \langle C,\sigma[\beta_s(p,\mathsf S,\varepsilon)\mapsto0]\rangle}
 \;\textsc{Scalar-Wait}
\end{gathered}
\]
\medskip
\begin{tabular}{@{}lll@{}}
\toprule
$a$ & $r$ & $\iota$\\
\midrule
$\kw{setflag}[p,q](v)$ & $p$ & $\kw{set}(q,\varepsilon)$\\
$\kw{waitflag}[p,q](v)$ & $q$ & $\kw{wait}(p,\varepsilon)$\\
$\kw{syncblockset}[p,q](v)$ & $p$ & $\kw{bset}(q,\varepsilon)$\\
$\kw{syncblockwait}[p,q](v)$，$q\ne\mathsf S$ & $q$ & $\kw{bwait}(p,\varepsilon)$\\
\bottomrule
\end{tabular}
\caption{操作与同步指令的发射；表格定义 $\mathsf{Enqueue}$}
\begin{itemize}
\item \textsc{Issue} 在发射时确定操作数与足迹；内存内容在操作完成时读取。
\item \textsc{Sync-Issue} 与 \textsc{Barrier-Issue} 只入队，控制立即继续；$\mathsf{ALL}$ 同时向本流各数据管道加入屏障。
\item \textsc{Scalar-Wait} 等待本流的块同步令牌，成功后只清除本流对应位；等待期间其他流及已入队操作仍可执行。
\end{itemize}
\end{figure}

\clearpage
\begin{figure}[!htbp]
\centering\small
\[
\begin{gathered}
\infer{Q_s(p)=U\cdot\kw{op}(f,\bar\nu,R,W)\cdot V\quad\mathsf{open}(U)
 \quad R\cup W\subseteq\mathrm{dom}(\mu)\cap\mathsf{Loc}_s
 \quad F_{\mathcal H}(s,p,f,\bar\nu,\mu|_R,\delta)}
 {\langle C,\sigma\rangle\xrightarrow{s}_{\mathcal H}
 \langle C,\sigma[\mu\mapsto\mu\oplus\delta,\ Q_s(p)\mapsto U\cdot V]\rangle}
 \;\textsc{Op-Done}\\[1em]
\infer{Q_s(p)=\kw{set}(q,\varepsilon)\cdot V\quad\varphi_s(p,q,\varepsilon)=0}
 {\langle C,\sigma\rangle\xrightarrow{s}_{\mathcal H}
 \langle C,\sigma[Q_s(p)\mapsto V,\ \varphi_s(p,q,\varepsilon)\mapsto1]\rangle}
 \;\textsc{Set-Done}\\[1em]
\infer{Q_s(q)=U\cdot\kw{wait}(p,\varepsilon)\cdot V\quad\mathsf{open}(U)
 \quad\varphi_s(p,q,\varepsilon)=1}
 {\langle C,\sigma\rangle\xrightarrow{s}_{\mathcal H}
 \langle C,\sigma[Q_s(q)\mapsto U\cdot V,\ \varphi_s(p,q,\varepsilon)\mapsto0]\rangle}
 \;\textsc{Wait-Done}\\[1em]
\infer{Q_s(p)=\kw{bar}\cdot V}
 {\langle C,\sigma\rangle\xrightarrow{s}_{\mathcal H}
 \langle C,\sigma[Q_s(p)\mapsto V]\rangle}
 \;\textsc{Barrier-Done}\\[1em]
\infer{Q_s(p)=\kw{bset}(q,\varepsilon)\cdot V\quad T=\mathsf{Targets}_{\mathcal H}(s,p,q)
 \quad\forall t\in T.\ \beta_t(p,q,\varepsilon)=0}
 {\langle C,\sigma\rangle\xrightarrow{s}_{\mathcal H}
 \langle C,\sigma[Q_s(p)\mapsto V,\ (\beta_t(p,q,\varepsilon)\mapsto1)_{t\in T}]\rangle}
 \;\textsc{Block-Set}\\[1em]
\infer{Q_s(q)=U\cdot\kw{bwait}(p,\varepsilon)\cdot V\quad\mathsf{open}(U)
 \quad\beta_s(p,q,\varepsilon)=1}
 {\langle C,\sigma\rangle\xrightarrow{s}_{\mathcal H}
 \langle C,\sigma[Q_s(q)\mapsto U\cdot V,\ \beta_s(p,q,\varepsilon)\mapsto0]\rangle}
 \;\textsc{Block-Wait}
\end{gathered}
\]
\caption{管道操作完成与令牌更新}
\begin{itemize}
\item $U,V$ 为队列序列；$\mathsf{open}(U)$ 表示 $U$ 中不含 \kw{wait}、\kw{bwait} 或 \kw{bar}。
队列分解指定本次移除的一个位置，即使存在相同队列项也不会一并删除。
\item \textsc{Op-Done} 要求足迹落在当前流可访问的已分配存储内（$R\cup W\subseteq\mathrm{dom}(\mu)\cap\mathsf{Loc}_s$），原子读取读区并提交更新；普通操作允许队内乱序完成（不必位于队首），以便在缺少 \kw{pipebarrier} 时暴露同管道依赖冲突。
\item \textsc{Set-Done}、\textsc{Block-Set} 与 \textsc{Barrier-Done} 必须在队首。
块通知要求所有接收位均已清零；各接收流独立消费自己的位。
\item 条件不满足时相应规则不启用；这些规则本身不将等待或数据竞争变成错误步骤。
\end{itemize}
\end{figure}

\subsection{终止、死锁与数据竞争}

给定可执行配置 $\langle P,\sigma\rangle$（其中 $P=\parallel_{s\in\kappa}s\{C_s\}$），定义三类核心检验性质：
\begin{itemize}
\item \textbf{正常终止（Terminated）}：对所有 $s\in\kappa$，$C_s\in\{\kw{return}(\bar v_s),\kw{end}\}$ 且对所有 $p\in\mathsf{Pipe}_{\mathsf{kind}(s)}$ 有 $Q_s(p)=\varepsilon$（基准配对检查另要求无残留未消费的同步令牌 $\varphi=0,\beta=0$）。
\item \textbf{死锁或卡住（Stuck / Deadlocked）}：$\langle P,\sigma\rangle$ 无法再执行任何一步 $\longrightarrow_{\mathcal H}$，且不满足正常终止条件。
\item \textbf{数据竞争（Data Race）}：存在位于队列中不同位置且均已打开（各自前缀满足 $\mathsf{open}(U)$）的两条在途操作项 $\kw{op}(f_1,\bar\nu_1,R_1,W_1)$ 与 $\kw{op}(f_2,\bar\nu_2,R_2,W_2)$（可在同一队列、同流不同队列或不同流中），满足
\[
(R_1\cap W_2)\cup(W_1\cap R_2)\cup(W_1\cap W_2)\ne\varnothing.
\]
\end{itemize}
\clearpage

\section{与 HIVM 的对应}

本节的 $\Delta_{\mathcal H}$、$\Sigma_{\mathcal H}$、$F_{\mathcal H}$ 和 $\mathsf{foot}_{\mathcal H}$ 只覆盖两个例子用到的运算名。$\kw{Copy}$、$\kw{MulAdd}$ 属于 $\mathsf{Operator}_{\mathsf{AIV}}$，$\kw{Fixpipe}$ 属于 $\mathsf{Operator}_{\mathsf{AIC}}$。复制关系 $F_{\mathrm{copy}}$ 按实际 view 操作数给出的对应单元传递当前内容，包括 $\mathsf{Uninit}$；足迹集合本身不决定对应顺序。

\begin{center}
\begin{tabularx}{\linewidth}{@{}llY@{}}
\toprule
HIVM & 核心形式 & 实例 \\
\midrule
\texttt{arith}、视图构造、\texttt{get\_block\_idx}
  & $\kw{scalar}[g]$
  & $\Delta$ 与 $\mathsf{eval}$ \\
\texttt{memref.alloc}、\texttt{memref.alloca}
  & $\kw{let}\ x=\kw{alloc}[m](\bar v)$
  & $\mathsf{alloc}_{\mathcal H}$，分配新鲜地址集并置 $\mathsf{Uninit}$ \\
\texttt{load}
  & $\kw{issue}[\mathsf{MTE2}](\kw{Copy},\ldots)$
  & $F_{\mathrm{copy}}$，GM 到 UB \\
\texttt{store}
  & $\kw{issue}[\mathsf{MTE3}](\kw{Copy},\ldots)$
  & $F_{\mathrm{copy}}$，UB 到 GM \\
\texttt{fixpipe}
  & $\kw{issue}[\mathsf{FIX}](\kw{Fixpipe},\ldots)$
  & $F_{\mathrm{copy}}$，L0C 到 GM \\
向量函数 \texttt{vf\_3}
  & $\kw{issue}[\mathsf{V}](\kw{MulAdd},\ldots)$
  & 三块输入都是值时按乘加写入目的地；含 $\mathsf{Uninit}$ 时目的地记为 $\mathsf{Uninit}$ \\
\texttt{set\_flag}、\texttt{wait\_flag}、\texttt{pipe\_barrier}
  & 同名核内命令
  & 本指令流的 $\varphi_s$ 与队列 \\
\texttt{sync\_block\_set}、\texttt{sync\_block\_wait}
  & $\kw{syncblockset}$、$\kw{syncblockwait}$
  & 按被通知指令流分开的 $\beta_t$ \\
\bottomrule
\end{tabularx}
\end{center}

\subsection{一个 AIV 上的乘加}

\texttt{debug-memory.mlir} 中的 \texttt{@mul\_add\_kernel} 使用 $\kappa_{0:1}=\{\mathsf{aiv}_1\}$，没有 \texttt{sync\_block}。三次装载分别使用事件 $0,1,2$。下面写出事件 $0$ 的一次装载，以及随后的乘加和两次写回；另外两次装载只更换事件和视图。源程序里由运行时决定的 \texttt{scf.if} 在进入该项之前取定分支，\texttt{set\_ctrl} 不进入语言。动态长度已经算进视图。
\[
\begin{aligned}
&\mathsf{aiv}_1\{
  \kw{setflag}[\mathsf{V},\mathsf{MTE2}](\varepsilon_0);\\
&\quad\kw{waitflag}[\mathsf{V},\mathsf{MTE2}](\varepsilon_0);
  \kw{issue}[\mathsf{MTE2}](\kw{Copy},\mathit{gm}_0,u_0);\\
&\quad\kw{setflag}[\mathsf{MTE2},\mathsf{V}](\varepsilon_0);
  \kw{waitflag}[\mathsf{MTE2},\mathsf{V}](\varepsilon_0);\\
&\quad\kw{issue}[\mathsf{V}](\kw{MulAdd},u_0,u_1,u_2,u_3);\\
&\quad\kw{setflag}[\mathsf{V},\mathsf{MTE3}](\varepsilon_0);
  \kw{waitflag}[\mathsf{V},\mathsf{MTE3}](\varepsilon_0);\\
&\quad  \kw{issue}[\mathsf{MTE3}](\kw{Copy},u_3,\mathit{gm}_3);
  \kw{issue}[\mathsf{MTE3}](\kw{Copy},u_3,\mathit{gm}_4);\\
&\quad\kw{pipebarrier}[\mathsf{ALL}];
  \kw{return}()\}.
\end{aligned}
\]
\texttt{vf\_3} 内部的 AVE 装载、乘法和加法都包在同一次 $\kw{MulAdd}$ 里。去掉 $\mathsf{MTE2}\to\mathsf{V}$ 这一对后，乘加项仍可在装载完成前打开并完成。它读到的是当时 UB 中的值。若那些单元仍是装载前的内容，写回 Global Memory 的就是这些内容的乘加结果。这是一次普通的完成步骤。

\subsection{AIC 与 AIV 的一次交接}

这次交接使用 $\kappa_{1:1}=\{\mathsf{aic},\mathsf{aiv}_1\}$。源文件是 \texttt{attention\_hivm\_lowered.mlir}。指令流 $\mathsf{aic}$ 对应函数 \texttt{\_attn\_fwd\_mix\_aic}。它用 \texttt{fixpipe} 写入 workspace，再执行 \texttt{sync\_block\_set}：核类型为 \texttt{CUBE}，管道为 \texttt{PIPE\_FIX} 与 \texttt{PIPE\_S}，flag 为 0。指令流 $\mathsf{aiv}_1$ 对应函数 \texttt{\_attn\_fwd\_mix\_aiv}。它先执行管道和 flag 相同、核类型为 \texttt{VECTOR} 的 \texttt{sync\_block\_wait}，再把 workspace 装入 UB。源指令中的核类型由 $\mathsf{kind}$ 给出，不另写进项。下面只保留这一次交接，不展开 \texttt{mmadL1}、\texttt{nd2nz} 和 FFTS 基址。L0C 中的值是 \texttt{fixpipe} 的输入。
\[
\begin{aligned}
&\mathsf{aic}\{
  \kw{issue}[\mathsf{FIX}](\kw{Fixpipe},\ell,w);
  \kw{syncblockset}[\mathsf{FIX},\mathsf{S}](\varepsilon_0);
  \kw{return}()\}\\
&\parallel
  \mathsf{aiv}_1\{
  \kw{syncblockwait}[\mathsf{FIX},\mathsf{S}](\varepsilon_0);
  \kw{issue}[\mathsf{MTE2}](\kw{Copy},w,u);
  \kw{return}()\}.
\end{aligned}
\]
$\kw{syncblockset}$ 要等同一条 \texttt{PIPE\_FIX} 队列里排在它前面的 \texttt{fixpipe} 离开，并且 $\mathsf{Targets}_{\mathcal H}(\mathsf{aic},\mathsf{FIX},\mathsf{S})$ 中每条流的对应位都为 $0$，才把这些位写成 $1$。在为本例配置相应的同步边后，$\kappa_{1:1}$ 里被通知的只有 $\mathsf{aiv}_1$。它的等待管道是 $\mathsf{S}$，因此在 $\beta_{\mathsf{aiv}_1}(\mathsf{FIX},\mathsf{S},\varepsilon_0)$ 变为 $1$ 之前不能发射后面的 \texttt{load}。等待完成时只清除 $\mathsf{aiv}_1$ 这一位。若换成 $\kappa_{1:2}$ 并在配置中提供到两条 AIV 流的同步边，则同一次 AIC 置位同时把 $\mathsf{aiv}_1$ 与 $\mathsf{aiv}_2$ 的对应位写成 $1$，两条 AIV 流各自清除自己的位。去掉 $\kw{syncblockwait}$ 后，$\mathsf{aiv}_1$ 可以先发射并完成装载，读到 \texttt{fixpipe} 写入之前 workspace 里的内容。

\appendix
\section{指令队列调查}
\label{app:queue-survey}

本附录记录 Ascend 910B 这一类分离核上的指令队列，不作为 \eqref{eq:hw} 的硬件参数。
深度取自 Ascend C 文档对 NPU 架构 2201 给出的指令队列深度，数据通路取自同一文档的搬运单元说明。
一个 AI Core 含一套 AIC 队列，以及 $\mathsf{aiv}_1$、$\mathsf{aiv}_2$ 各一套 AIV 队列。
深度为 $0$ 表示该核没有这条执行单元。

\begin{center}
\small
\begin{tabularx}{\linewidth}{@{}llYcYc@{}}
\toprule
队列 & 执行单元 & AIC 上的数据通路 & AIC 深度
  & AIV 上的数据通路 & AIV 深度 \\
\midrule
$\mathsf{MTE2}$ & 搬入 DMA
  & $\mathsf{GM}\to\mathsf{L1}$，也可 $\mathsf{GM}\to\mathsf{L0A}/\mathsf{L0B}$ & $32$
  & $\mathsf{GM}\to\mathsf{UB}$ & $16$ \\
$\mathsf{MTE1}$ & L1 往 Cube 入口搬
  & $\mathsf{L1}\to\mathsf{L0A}$、$\mathsf{L1}\to\mathsf{L0B}$、$\mathsf{L1}\to\mathsf{BiasTable}$ & $32$
  & 无此队列 & $0$ \\
$\mathsf{M}$ & Cube
  & 读 $\mathsf{L0A}$、$\mathsf{L0B}$，读写 $\mathsf{L0C}$ & $16$
  & 无此队列 & $0$ \\
$\mathsf{FIX}$ & FixPipe
  & $\mathsf{L0C}\to\mathsf{GM}$、$\mathsf{L0C}\to\mathsf{L1}$ & $32$
  & 无此队列 & $0$ \\
$\mathsf{MTE3}$ & 搬出 DMA
  & $\mathsf{L1}\to\mathsf{GM}$ & $32$
  & $\mathsf{UB}\to\mathsf{GM}$ & $16$ \\
$\mathsf{V}$ & Vector
  & 无此队列 & $0$
  & 读 $\mathsf{UB}$，写 $\mathsf{UB}$ & $64$ \\
\bottomrule
\end{tabularx}
\end{center}

\noindent
Scalar 是每条流的发行者，上述深度表没有给出 Scalar 的队列深度。$\mathsf{ALL}$ 不是一条指令队列。

\end{document}
